\chapter{Benchmark strategies}

The first strategies we implemented are the three benchmarks of the book
\parencite[Ch.~3]{lihoi}: buy and hold, the constant rebalanced portfolio and the best
stock. They are the reference point, since an algorithm that cannot beat buy and hold is not worth using due to transaction costs.
Implementing these benchmarks first was also useful to test the simulation with strategies whose results we can check by
hand. In all of them the weights are equal at the beginning, $b=(1/m,\dots,1/m)$,
unless said otherwise.

\section{Buy and hold}

Buy and hold (BAH) is the simplest strategy. We choose a portfolio $b_1$ on the first
day and never trade again. Its final wealth is
\[
S_n(\mathrm{BAH}) = b_1\cdot\bigodot_{t=1}^{n} x_t .
\]
The formula looks odd at first, so we read it piece by piece. The symbol $\bigodot$ is
the element-wise product of the vectors $x_1,\dots,x_n$, so the $i$-th component of
$\bigodot_t x_t$ is $\prod_t x_{t,i}$. That is the factor by which the price of asset
$i$ has been multiplied from the first to the last day, for example 2 if it has
doubled or 0.5 if it has lost half of its value. The dot product with $b_1$ gives the
average of these factors weighted by the initial weights. It makes sense, because the
money we put in asset $i$ has grown exactly as much as the price of that asset.

Even if we never trade, the weights of the portfolio do change. An asset that goes up
becomes a bigger part of the total, and one that goes down becomes smaller. On day
$t+1$ the weights are
\[
b_{t+1,i}=\frac{b_{t,i}\,x_{t,i}}{b_t\cdot x_t},
\]
which is the old weight multiplied by the return of the asset and divided by the
return of the whole portfolio. We need this expression in the code, because the
simulation asks every strategy for its portfolio each day, including BAH.

\section{Constant rebalanced portfolio}

The constant rebalanced portfolio (CRP) \parencite{cover91} does the opposite. It fixes
a portfolio $b$ and goes back to it every day, so its final wealth is
\[
S_n(\mathrm{CRP}(b)) = \prod_{t=1}^{n} b\cdot x_t .
\]
Each factor $b\cdot x_t$ is how much the wealth grows on day $t$ (the average of the
price relatives of that day weighted by $b$), and the wealth after $n$ days is the
product of the daily factors. The difference with BAH is where the product is placed.
In BAH it is inside the dot product, so each asset grows on its own. In CRP it is
outside, because every day we start again from the same weights.

After the prices move, the weights are no longer $b$. To go back to $b$ the strategy
has to sell a bit of what has gone up and buy a bit of what has gone down. So CRP is a
contrarian strategy, and it trades every day, which we ignore for now because the
simulation has no transaction costs. The reason why it is interesting is that it can
take advantage of the fluctuations. If one asset goes up and then comes back down
while another does the opposite, BAH ends where it started, but CRP has been selling
high and buying low in between. In general, the growth rate of a CRP is never lower
than the weighted average of the growth rates of its assets, and it is higher the more
the assets move differently from each other.

\section{Best stock and BCRP}

The best stock puts all the wealth in the asset that grew the most over the whole
period,
\[
S_n(\mathrm{Best\ Stock}) = \max_{i=1,\dots,m}\ \prod_{t=1}^{n} x_{t,i}.
\]
It is a BAH whose initial portfolio is concentrated on the best asset. The problem is
that we only know which asset is the best once the period is over. This means that it
cannot be run in real time, because it uses information from the future. It is not a
competitor but a ceiling to compare with.

There is a similar ceiling for CRP, the best constant rebalanced portfolio (BCRP),
which is the CRP with the weights that would have been the best in hindsight,
\[
b^\ast=\arg\max_{b\in\Delta_m}\prod_{t=1}^{n} b\cdot x_t .
\]
We do not compute it in this chapter, but it will be the target of the algorithms of
the next one. The three ceilings are ordered:
\[
S_n(\mathrm{BAH}(b_1))\le S_n(\mathrm{Best\ Stock})\le S_n(\mathrm{BCRP}).
\]
The first inequality holds because BAH gives an average of the growth factors of the
assets, and an average cannot be larger than the largest of them. The second holds
because the best stock is also a CRP, the one with all the weight on a single asset
(it never has to rebalance), so it is one of the portfolios among which the BCRP is
chosen. The ordering only applies to BAH. A particular CRP, like the uniform one, can
end below or above the best stock.

\section{Implementation}

All strategies share one interface. A strategy receives the price relatives observed
up to $t-1$ and returns the portfolio $b_t$. On the first day, when there is no
history, the strategy returns its initial portfolio, which is the equal-weight one
unless the strategy sets another. Every portfolio returned is validated before it is
used: it must have $m$ components, none negative, and add up to one. The simulation
asks for $b_t$ with only $x_1,\dots,x_{t-1}$ in hand, multiplies it by $x_t$ and
updates the wealth with $S_t=S_{t-1}\,(b_t\cdot x_t)$, $S_0=1$.

The BAH class computes the drifted portfolio from the history it receives. The CRP
class just returns its fixed $b$. The best stock is written as a BAH whose initial
portfolio is chosen in the constructor from the complete history. This is the only
place in the code where a strategy sees future data, which is why it is kept apart and
never used as a comparison with online algorithms.

We checked the three classes on a series of 500 days and four assets simulated
with random price relatives. The final wealth that the simulation returns for each
strategy matches the formulas of this chapter, differing by less than $10^{-13}$.
